% Volume V: Admissibility and Continuation
% Part IX. A General Theory of Institutional State Change
% Chapter 50: The Admissibility Boundary
% Spine: proof-spines/ch050-spine.md

\chapter{The Admissibility Boundary}
\label{ch:ch050}

Let \(\Omega\) be the space of conceivable claims and \(\mathcal D_c\subseteq\Omega\) the claims admissible in context \(c\). The \textbf{admissibility boundary} is then
\[
\partial\mathcal D_c,
\]
\Definition\ where \Omega is the claim-space and \(c\) the forum context.
which contains claims whose status is contested, incomplete, or dependent on further evidence.

Admissibility is therefore not identical with truth. It is a forum-specific judgment about which claims may presently bear consequence under given rules and capacities of interpretation. Much institutional conflict concerns motion of the boundary rather than direct disagreement about truth: actors fight over whether a claim counts as ready, legible, or procedurally fit for uptake at all.
